\reviewercomment{6. Mathematical model

The gLV model successfully reproduces the qualitative transition in coalescence outcomes and serves as a useful minimal framework. However, it is limited to purely competitive interactions and does not explicitly capture environmentally mediated effects such as pH modification.

We therefore suggest framing the model more explicitly as a phenomenological rather than mechanistic description of the system.

In addition, while the robustness analysis across interaction coefficient distributions is appreciated, a brief discussion of the biological interpretation of these distributions would be helpful. In particular, clarifying how the mean interaction strength parameter ($\mu$) relates to underlying ecological processes would improve interpretability.}

\responselabel We agree that the gLV model should be presented as a phenomenological competitive baseline, not as a complete mechanistic model of the experimental system. We have therefore revised the model description to state explicitly that the gLV framework explores how the statistical properties of effective interaction coefficients shape coalescence outcomes, while omitting explicit environmental state variables such as pH, metabolite concentrations, and carrying-capacity differences. In the model, $\alpha_{ij}$ is now described as an effective per-capita inhibitory term for the effect of species $j$ on species $i$; it can absorb the net steady-state consequences of direct competition and indirect inhibition, including pH-mediated inhibition, but it does not identify the biochemical mechanism and does not dynamically model pH itself. Because we restrict $\alpha_{ij} \geq 0$, we also now state that the simulations represent competitive or growth-inhibiting interactions only, and that facilitative interactions such as cross-feeding are not represented.

We also revised Supplementary Methods to define $\alpha_{ij}$ as an effective coarse-grained inhibitory coefficient and to state that the baseline $\alpha_{ij}\geq 0$ model captures competitive growth inhibition rather than facilitation. The uniform, Gaussian, and Gamma coefficient distributions are phenomenological ensembles of heterogeneous effective competitive effects, not fitted biochemical distributions. Biologically, $\mu$ should be read as a minimal phenomenological axis for interaction-mediated assembly structure, not as a direct measurement of nutrient supply or a fitted biochemical parameter. In the baseline model, increasing $\mu$ bundles two effects: more pairwise coefficients become appreciably different from the no-interaction limit at zero, and the coefficients are sampled from a broader, more heterogeneous range. The first effect moves the system away from near-neutral coexistence toward appreciable growth inhibition; the second creates stronger mismatches between species that assembled together and outsiders, allowing assembly to filter communities into internally compatible groups that resist invasion together. We decoupled these two axes in Supplementary Note~3. The mean-vs-variance sweep shows that coefficient mean alone is not sufficient; both average inhibitory effect and coefficient heterogeneity contribute to the Dominance regime.

This reframing preserves the main use of the model while making its limits explicit. The competition-only gLV model shows that competitive assembly structure is sufficient to generate frequent Dominance and positive within-parental-community selection correlation, but it does not claim that competition is the only mechanism operating in the experiments. Indeed, the revised Results and Discussion now emphasize that pH modification by dominant taxa is one possible contributor to winner identity in Nutr$+$: in the strict acidic--alkaline subset used in Extended Data Fig.~7, the acidic parental community dominated 29/32 Nutr$+$ events (90.6\%), whereas the Base signal was weaker (23/41 events, 56.1\%). Thus, the model is used to identify a sufficient coarse-grained mechanism for community-level selection, whereas the experimental interpretation allows environmental mediation to contribute to additional features of the data without treating pH as a complete explanation for Nutr$+$ Dominance frequency.

\mschange{We revised the Results to call the gLV model a phenomenological framework rather than a mechanistic model of pH modification, metabolic cross-feeding, or carrying-capacity variation. We revised the Supplementary Methods to define $\alpha_{ij}$ as an effective inhibitory coefficient, to state that $\alpha_{ij} \geq 0$ makes the model competitive-only, and to define $\mu$ as a coarse-grained interaction-intensity axis that scales the baseline coefficient distribution. We also revised the Discussion to state that the mapping between nutrient concentration and $\mu$ is approximate and reflects net interaction intensity measured through invasion resistance, not a unique biochemical mechanism. No new response figure was added because the concern is addressed by textual clarification and by existing main-text and supplementary analyses.}
